\chapter*{Reflexión final}
\addcontentsline{toc}{chapter}{Reflexión final}

Este trabajo me permitió explorar no solo las capacidades técnicas de diferentes modelos, sino también las implicancias prácticas de aplicar herramientas de inteligencia artificial en un contexto clínico real y sensible. A partir de esta experiencia, considero que hay varias direcciones que podrían ser exploradas en futuras investigaciones.

Algo que me gustaría seguir profundizando es cómo lograr que modelos como GPT-4o se acerquen un poco más a criterios clínicos reales. Cambiar el prompt o las imágenes de ejemplo puede modificar totalmente lo que el modelo responde, y eso muestra tanto su potencial como sus límites. Quizás una forma de mejorar esto sea combinar este tipo de modelos con otros sistemas que tomen decisiones más controladas, como reglas clínicas o redes entrenadas específicamente para una patología.

También sería valioso desarrollar marcos de evaluación específicos para medir si estas respuestas generadas son verdaderamente útiles para profesionales de la salud, más allá de métricas numéricas. Porque a veces acierta la etiqueta, pero no necesariamente da una explicación útil o coherente desde lo clínico.

Finalmente, aunque las soluciones de bajo esfuerzo como GPT-4o son prometedoras, siguen siendo herramientas auxiliares. Por ahora, su uso en salud requiere acompañamiento, validación y responsabilidad humana. Pero con una integración cuidadosa, podrían ocupar un lugar útil en sistemas híbridos de apoyo al diagnóstico, especialmente en contextos donde los recursos son escasos y el acceso temprano a la atención puede marcar una diferencia significativa.